\section{Ausblick}
\label{chap:ausblick}

Für die Überführung der Ergebnisse in eine serienfähige Schaltanlage sind als nächste Schritte insbesondere Validierungen unter realitätsnahen Randbedingungen erforderlich. Dazu zählt zunächst die konstruktive Umsetzung der empfohlenen Anordnung in einem Prototypenaufbau der Neuentwicklung. In diesem Schritt sollte die Messstabilität über den vorgesehenen Strombereich erneut verifiziert werden, wobei die Phase L2 weiterhin als kritischer Indikator für die Fremdfeldrobustheit zu betrachten ist. Ergänzend ist eine Robustheitsuntersuchung gegenüber Einbauvariationen sinnvoll, da in der Serienfertigung Toleranzen in Abstand, Ausrichtung und Leiterführung unvermeidbar sind und sich auf die Feldkopplung auswirken können. Ziel ist es, ein Prozessfenster abzuleiten und montagegerechte Designregeln (z.\,B. Mindestabstände, definierte Einbaulagen) festzuschreiben.

Darüber hinaus sollte die Variantendeckung geprüft werden. Es ist zu klären, ob das empfohlene Basiskonzept über unterschiedliche Leistungsklassen und Geometrievarianten hinweg stabil bleibt oder ob in einzelnen Konfigurationen eine modulare Ergänzung durch FFP erforderlich wird. Ergänzend ist ein Vergleich von Messstromwandlern mit unterschiedlichen Nennbürden sowie der Einfluss unterschiedlicher Bürdenkonfigurationen zu untersuchen. Im Fokus stehen definierte Sekundärwiderstände und die resultierende Sekundärlast, um deren Auswirkung auf Übersetzungsmessabweichung und Stabilität unter praxisnahen Bedingungen zu quantifizieren und daraus gegebenenfalls eine empfohlene Bürdenauslegung für den Zielanwendungsfall abzuleiten.



Der Hochstromprüfstand wurde im Rahmen dieser Arbeit bereits weiter optimiert. Perspektivisch sind Erweiterungen sinnvoll, um zusätzliche Versuche unter seriennahen Betriebsbedingungen durchführen zu können. Dazu zählen insbesondere Erwärmungsprüfungen sowie eine automatisierte Wandlerprüfung über definierte Prüfsequenzen. Hierfür ist eine Erweiterung der Steuerung und der Prüfhardware erforderlich, sodass im Prüffeld auch komplette Baugruppen oder Anlagen angeschlossen und als Gesamtsystem über standardisierte Prüfabläufe geprüft werden können. Ergänzend ist eine automatisierte Erfassung, Auswertung und Berichtserstellung sinnvoll, um die Prüfergebnisse reproduzierbar zu dokumentieren und kundenfähige Prüfberichte bereitzustellen.